%! Author = leste
%! Date = 2026-09-14

% Preamble
\documentclass[%
 reprint,
%superscriptaddress,
%groupedaddress,
%unsortedaddress,
%runinaddress,
%frontmatterverbose,
%preprint,
%preprintnumbers,
nofootinbib,
nobibnotes,
%bibnotes,
 amsmath,amssymb,
 aps,
%pra,
%prb,
%rmp,
%prstab,
%prstper,
floatfix,
]{revtex4-2}

% Packages
\usepackage{graphicx}% Include figure files
\usepackage{dcolumn}% Align table columns on decimal point
\usepackage{bm}% bold math
\usepackage{siunitx}
\usepackage{booktabs}
\usepackage{derivative}
\usepackage{amsmath}
\usepackage{mathtools}

%\usepackage[giveninits=true]{biblatex}
%\addbibresource{references.bib}

% Recommended options for physics papers:
\sisetup{
  separate-uncertainty = true, % Prints as (value ± uncertainty) unit instead of value(unc) unit
  multi-part-units = single    % Ensures units are not repeated, e.g., (1.23 ± 0.04) m instead of 1.23 m ± 0.04 m
}


\newcommand{\UVic}{\affiliation{Department of Physics and Astronomy, University of Victoria, Canada}}
\newcommand{\dif}{\mathop{}\!\mathrm{d}}

% Document
\begin{document}

% \preprint{APS/123-QED}

\title{Experiment 33: Electromotive Response To Moving Magnet Sources}

\author{Mark Lester}\UVic

\date{\today}% It is always \today, today,
             %  but any date may be explicitly specified

\begin{abstract}

%Presents the major aspects of the investigation and includes:
%1) the overall purpose of the lab;
%2) the basic design of the
%investigation;
%3) major findings or trends found;
%4) a summary of interpretations and conclusions.

    This laboratory experiment aims to study the behaviour of EMF generated by cylindrical magnets falling through
    conductive coils of wire.
    The data collected will be compared with magnetic flux equations, Faraday's Law, Lenz' Law, in order to establish a
    basic framework for magnetostatics.
    Our results show that while magnetic field strength is sensitive to local sources, general agreement with
    calculated EMF curves is possible given rigid parameters and precise measurements.
    It was found that a mechanical-electrical energy transfer was represented as an EMF, and that the current
    induced in the coils from a moving magnet produced a back-EMF that opposed the direction of the magnet's field
    lines, which in turn marginally slowed the descent.

\end{abstract}


%\keywords{Suggested keywords}%Use showkeys class option if keyword
 %display desired
\maketitle

 %\tableofcontents
\section{Introduction}


%Presents a clear learning context and includes:
%1) the general scientific concept with the necessary background;
%2) primary goals and the specific scientific concept to be investigated;
%2) specific theory, experimental design and equations needed with fully defined quantities;
%3) any assumptions made or required;
%4) hypothesis to be investigated;
%5) references to supporting sources.

Electromotive Response, or EMF, is responsible for the mechanical-electrical energy transfer across generators,
motors, and transformers.
The objective of this study is to demonstrate this exchange in the form of magnetic dipoles accelerating
conductive coils.

\begin{equation}\label{eq:faraday}
    \epsilon_1 = -\frac{\dif \phi}{\dif t}
\end{equation}

In Eq.~\eqref{eq:faraday}, Faraday's Law describes that a change in flux results in an EMF\@.

By forming the differential equation that relates flux to the magnetic field, we can show that,

\begin{equation}\label{eq:flux}
    d \phi = \frac{\mu_0 p_m}{2 \pi} \frac{z^2 - r^2 / 2}{(r^2 + z^2)^{5/2}} (2 \pi r dr)
\end{equation}

where $d \phi$ is the flux element, $p_m$ is the magnetic dipole moment, $z$ is the axial distance away from the
coil, $\mu_0$ is the permeability of free space, and $r$ is the radius of the loop.

Further, we can derive the EMF by integrating across the time element.

\begin{equation}\label{eq:emf}
    \epsilon = -\frac{3}{2} N \mu_0 p_m r_0^2 \frac{z}{(r_0^2 + z^2)^{5/2}}v
\end{equation}

Where $\epsilon$ is the EMF, $N$ is the number of coil turns, $r_0^2$ is the coil radius, and $v$ is the speed of
the magnet moving through the coil.

Additionally, for sufficiently large axial distances, Eq. ~\eqref{eq:Bz} applies.
This allows us to determine the magnetic dipole moment from $B_z$ and distance measurements.

\begin{equation}\label{eq:Bz}
    B_z \approx p_m \frac{\mu_0}{2 \pi} \frac{1}{z^3}
\end{equation}


%Ref equation: Eqs.~\eqref{eq:norm_damp}, ~\eqref{eq:crit_damp}, and \eqref{eq:osc_damp}, respectively.
\section{Apparatus}

%Equipment: Vernier LabPro sensor interface, LabPro software, 4000 turn inductor, electrolytic capacitor, decade
%resistance box...

Four coils, each having $N=20$, $r_0 = 0.0442$m, cylindrical magnets of 2.5cm and 8.5cm lengths, dual input
amplifiers, DC Gaussmeter Model GM-1-ST, LabPro sensor interface, LabPro software, vertical support rod.

    \begin{figure}[htb]
        \centering
        \includegraphics[width=\linewidth,height=0.35\textheight,keepaspectratio]{figures/apparatus}
        \caption{Vertical rod holding the coils}
        \label{fig:diag}
    \end{figure}

In Figure ~\ref{fig:diag}, the vertical rod anchored the coils 250mm apart, with a weighted guiding string threaded
through the centre.
The coils were connected to the amplifiers of gain $300 \pm 1\%$.
\section{Procedure}\label{sec:procedure}

%Presents the apparatus and procedure used and includes:
%1) a list of the apparatus;
%2) an apparatus diagram; 3) the technical details of the investigation in the procedure used to acquire data.
%Written in the past tense with specific terminology, this should allow others to replicate the experiment.

The procedure below followed the outline in the PHYS 326 course lab book ~\cite{phys326}.

Two methods were used to find the dipole moment of each magnet: Method 1 used a Gauss meter to measure the magnetic
field strength
and plotting the far-field behaviour, and method 2 used the EMF curve from dropping the magnets through successive
coils.

A ruler was used to measure the axial distance $z$ whilst placing the Gauss meter at regular intervals away from the
magnet.
$B_z$ and their corresponding distance values were recorded in a table and incremented by 5mm until the magnetic
field strength became too weak to measure.
This step was done for both the short and long magnets.

\begin{figure}[htb]
        \centering
        \includegraphics[width=\linewidth,height=0.35\textheight,keepaspectratio]{figures/measure}
        \caption{Method 1: The magnet was placed along the ruler. $B_z$ and corresponding distance measurements were taken.}
        \label{fig:meas}
    \end{figure}

The slope of the line where it was approximately linear was used to derive the magnetic dipole moment using Eq. ~\eqref{eq:Bz}.

As for method 2, a magnet was carefully attached to the top of the string and left to fall (with North polarity
facing down) from rest.
With the coils hooked up to dual input amplifiers, this resulted in EMF which was sampled by the LabPro sensor
hardware.
This data was recorded in LabPro at the highest sampling rate available, and EMF curves for each coil were plotted
against time.
The EMF curves were de-amplified by the gain rating used to reflect the actual response.
This was repeated for both the short and long magnets.

\clearpage


%The experimental apparatus and procedure followed the standard method outlined in the University of Victoria PHYS 326 laboratory manual~\cite{phys326}.

%Table~\ref{tab:params} provides a list of measurements taken.
%Uncertainty was calculated using...


%\begin{table}[htbp]
%\caption{Experimental Parameters}
%\begin{tabular}{lrrr}
%\toprule
%{Quantity} & {Reading} \\
%\midrule
%$R_L$, Inductor resistance & \qty{109.4(1.3)}{\ohm} \\
%L, Inductance & \qty{1.044(0.072)}{\henry} \\
%C, Capacitance & \qty{10.07(0.21)}{\micro\farad}\\
%$R_{FG}$, function generator resistance & \qty{50.00(0.25)}{\ohm} \\
%$R_{res}$, normal & \qty{1133}{\ohm} \\
%$R_{res}$, critical & \qty{485}{\ohm} \\
%$R_{res}$, oscillatory & \qty{20}{\ohm} \\
%\bottomrule
%\end{tabular}\label{tab:params}
%\end{table}
\section{Results}\label{sec:results}

%Presents the results of the procedure and includes:
%1) all data collected in tabular form if possible;
%2) all required graphs;
%3) all results in a logical sequence;
%4) objective observations;
%Note: all figures, tables and graphs required should include
%titles, labels, units and legends where appropriate.

%Present processed data, not raw data.

\begin{figure}[htb]
    \centering
    \includegraphics[width=\linewidth]{figures/magnet_strength}
    \caption{Magnetic field strength of the long and short magnets as a function of inverse cube distance. }
    \label{fig:method1}
\end{figure}

In Figure ~\ref{fig:method1}, the magnetic field strength was plotted against inverse cube distance in order to
approximate the value of the dipole moment.
As the dipole moment equation is only a far field approximation, the theoretical linear relationship corresponds to
the data collected the furthest away from the magnet.

\begin{figure}[htb]
    \centering
    \includegraphics[width=\linewidth]{figures/short_fall}
    \caption{EMF induced by releasing the magnetic from rest.}\label{fig:short_fall}
\end{figure}

\begin{figure}[htb]
    \centering
    \includegraphics[width=\linewidth]{figures/long_fall}
    \caption{EMF induced by releasing the long magnet from rest.}\label{fig:long_fall}
\end{figure}

In Figures ~\ref{fig:short_fall} and ~\ref{fig:long_fall}, the EMF is induced from the flux of the magnetic field
lines changing with time as the magnet falls.
The voltage peak-peak amplitude increases with each successive coil the magnet falls through.

By performing a Riemann sum over the positive and negative values, we can resolve the flux through each coil.
Table ~\ref{tab:riemann1} and ~\ref{tab:riemann2} shows the calculated result from the falling short magnet.
Each coil has close agreement between their positive and negative phases, resulting in approximately equal and
opposite total flux.

\begin{table}[htb]
    \caption{Calculated Flux Values - Positive Phases}
    \centering
\begin{tabular}{lr}
    \toprule
Flux Component & Quantity {[}Wb{]} \\
    \midrule
fluxA          & -1.117e-04        \\
fluxB          & -1.081e-04        \\
fluxC          & -1.133e-04        \\
fluxD          & -1.104e-04        \\
flux total     & -2.218e-05        \\
    \bottomrule
\end{tabular}\label{tab:riemann1}
\end{table}

\begin{table}[htb]
    \caption{Calculated Flux Values - Negative Phases}
    \centering
\begin{tabular}{lr}
    \toprule
Flux Component & Quantity [Wb] \\
    \midrule
fluxA          & 1.106e-04        \\
fluxB          & 1.073e-04        \\
fluxC          & 1.126e-04        \\
fluxD          & 1.095e-04        \\
flux total     & 2.200e-05        \\
    \bottomrule
\end{tabular}\label{tab:riemann2}
\end{table}



\begin{figure}[htb]
    \centering
    \includegraphics[width=\linewidth]{figures/position_short}
    \caption{Z position as a function of time.}\label{fig:position_short}
\end{figure}

\begin{figure}[htb]
    \centering
    \includegraphics[width=\linewidth]{figures/position_long}
    \caption{Z position as a function of time.}\label{fig:position_long}
\end{figure}

In Figures ~\ref{fig:position_short} and ~\ref{fig:position_long}, a quadratic fit is applied to the position data
at the points where the magnet passes through the coils.
By extracting the value from the quadratic fit line, we determine that the acceleration due to gravity was
$\qty{9.41}{\meter\per\second^2}$ and $\qty{9.48}{\meter\per\second^2}$ for the short and long magnets,
respectively.
To derive an expression for speed, $v(t)$, we take the derivative of $z(t)$:

\begin{equation}
    v(t) =
    \begin{cases*}
        9.41t + 3.54 & short magnet \\
        9.48t + 3.46 & long magnet
    \end{cases*}
    \label{eq:velocity_cases}
\end{equation}

\begin{figure}[tb]
    \centering
    \includegraphics[width=\linewidth]{figures/short_coilB}
    \caption{Experimental and calculated EMF curve for short magnet.}\label{fig:short_coilB}
\end{figure}


\begin{figure}[tb]
    \centering
    \includegraphics[width=\linewidth]{figures/long_coilB}
    \caption{Experimental and calculated EMF curve for long magnet.}\label{fig:long_coilB}
\end{figure}

In Figures ~\ref{fig:short_coilB} and ~\ref{fig:long_coilB}, we compare the experimental and calculated EMF for
coil.

The calculated EMF response from the long magnet deviates more significantly compared with the short magnet.


%\clearpage


%\textcolor{red}{
%    Uncertainties in experimental values were calculated using the aid of the \textit{uncertainties} python package.
%}

\section{Discussion}\label{sec:discussion}

%Presents a clear statement and insightful interpretation of the results, accurately relates results to the hypothesis, and discusses any limitations or sources of error in the experiment.
%Questions are answered completely and correctly.


The collected data showed that dropping magnets from rest through coils induced an EMF due to a time-varying flux.
The amplitude of each successive coil increased in magnitude as the magnet increased in speed, and the phase switched polarity where the magnets passed through each centre.

By using our derived magnetic dipole moment values from method 1, we calculated a theoretical EMF curve.
This was overlaid on our coil B data as shown in Figures ~\ref{fig:short_coilB} and ~\ref{fig:long_coilB}.
Though the theoretical signal was greater in amplitude compared with our experimental result, the overall shape and
behaviour was consistent with the theory.

\bigskip

The attenuation seen in the experimental result can be attributed to inefficiencies in exchanging the mechanical
energy of the falling magnet with the electrical energy.
Additionally, other magnetic field lines produced from local sources, such as nearby computers and cables, also
affected the EMF
generated by
the
setup.
Further reduction in EMF can be explained by the induced current transferring its energy into Joule heating.

\bigskip

Our analysis on the quadratic interpolation of position showed that the magnets' acceleration was marginally less
than the acceleration due to gravity, $9.81 m/s^2$.
This suggests there is an opposing force to gravity which led to experimental accelerations of $9.41 m/s^2$ and
$9.48 m/s^2$.
We can attribute this to Lenz' Law, where magnetic field generated by the falling magnets resists the direction of
motion.
Instead of using $9.81 m/s^2$ for calculations, the experimental accelerations were used because it
reflected the physical reality more closely.


Due to the geometry of the magnets, the magnetic field lines did not adhere to theoretical $B_z$ when measuring the
field strength near the magnet.
This behaviour changed and $B_z$ began to scale approximately linearly with $z^{-3}$ once the Gauss sensor was far
enough away.
For Figure ~\ref{fig:method1}, the short magnet showed a linear trend for $B$ versus $z^{-3}$, whereas for the majority
of the data points, the
long magnet did not.
The most linear sections of the data were taken in order to find the far field approximation for the magnetic dipole
moment.
This was a problem for the long magnet, as it did not show much linear dependence until further away from the
magnet.
If farther measurements had been taken, the long magnet would have shown a greater linear dependence on $z^{-3}$ and
produced a shallower slope.
This would have been reflected as closer fit to the data for the calculated EMF curve seen in Figure ~\ref{fig:long_coilB}.

\bigskip

Table~\ref{tab:riemann1} and~\ref{tab:riemann2} show the Riemann sum of positive and negative values for
flux across each coil for the short magnet.
This demonstrated that the flux across both phases were approximately equal and opposite for each coil.
The sum total for positive and negative phases was -2.218e-5 Wb and 2.200e-5 Wb, respectively.
This result agrees with the theory since the path of the magnet on either side of each coil is symmetric about the z
axis, and flux is found from Faraday's law, which is dependent on $z$ position.

According to ~\eqref{eq:flux}, there is an upper and lower limit to maximising the magnetic flux.
By taking the numerator, $z^2 - r^2 / 2$, we can determine that if the radius of the coil exceeds $\sqrt{2}z$, then
the flux element starts to become negative.
Therefore, to maximise the net flux as a result of coil radius, we would ideally have a coil with dynamic radius
that would expand or contract to $\sqrt{2}z$.
The lower limit is the radius of the magnet itself.
However, the radial element increases at the same rate as $B_z$ decreases, thus cancelling the scaling effect.
This also explains why the calculated flux is symmetrical for both positive and negative phases despite the
acceleration of the
magnet.

\section{Conclusion}\label{sec:conclusion}

%Presents a clear, concise, focused, insightful and correct connection between the results and the investigated
%concepts.
%Explains why the findings are important, what they show, how they contribute to the broader understanding of the topic and recommendations for future research or improvements to
%experimental design.
%The section brings closure to the report.

It was discovered that a moving magnet produces a time-varying flux and thus an EMF through
multi-turn
coils.
The resulting induced current also form its own magnetic field which opposed the direction of the field lines of the
magnet, agreeing with Lenz' Law.
Method 1 yielded magnetic dipoles of \qty{0.56}{\ampere\meter^2} and \qty{2.05}{\ampere\meter^2} for the short and
long magnets,
respectively.
Though these values resulted in overshooting the calculated amplitude of the EMF versus the experimental signal, the
theory showed an agreement with the overall shape of the response curve.

Deviations from the calculated signal were attributed to other local magnetic field sources, geometric limitations, and
Joule heating.
For better results, the experiment should be performed in a container which shields itself from other magnetic fields.

\clearpage            % Flushes two-column floats and starts a fresh page
\onecolumngrid        % Switches REVTeX to full page width

\appendix
\section{Sample Calculations}\label{sec:sample-calculations}

% sample calcs here
Magnetic dipole moment calculation:

\begin{align}
    p_m &= (\text{slope}) \frac{2 \pi}{\mu_0} \\
    &= (1.12e-7) \frac{2 \pi}{1.25663706e-6} \qty{}{\ampere\meter^2} \\
    &= \qty{2.05}{\ampere\meter^2}
\end{align}

FLux calculation by Riemann sum per phase.
Calculation was done using np.sum in python.
This is an example for the positive phase for coil A:

\begin{align}
    \phi &= \sum_{t_1}^{t_2} \frac{1}{N}\epsilon(t_n) \Delta t \\
    &= \sum_{0}^{0.0656} \frac{1}{20}\epsilon(t_n) (0.0004) \\
    &= -1.117e-04 \, \text{Wb}
\end{align}

Calculated EMF curve.
The calculation was done using vectorisation of a numpy array in python, but this is the first calculated data point:

\begin{align}
    \epsilon &= \frac{-3}{2} N \mu_0 p_m r_0^2 \frac{z}{(r_0^2 + z^2)^{5/2}} v \\
    &= \frac{-3}{2} (20) (1.2566e-6) (2.054) (0.0442)^2 \frac{(-0.627)}{((0.0442)^2 + (0.0442)^2)^{5/2}} (-0.332) \qty{}{\volt} \\
    &= \qty{-3.27e-7}{\volt}
\end{align}

\section{Source code and raw data}\label{sec:source-code-and-raw-data}

The source code and raw data may be found in the public git repository:

\url{https://github.com/trolley-lorry/phys_bsc_data/tree/2535c1549e291a68488e0db284f57de73a7c8ead/phys326/lab33}.


\bibliography{references}
\bibliographystyle{apsrev4-2}
%\bibliographystyle{plain}



\end{document}